% Self-contained report; local Forbes/Lucida report style.
% Compile twice with /Library/TeX/texbin/pdflatex.
% Adapted from the 7 September comparison in the parent research checkout.
% Standalone core status refreshed on 9 September 2026; tooling assessment remains dated 7 September.
\documentclass[12pt,leqno]{article}
\usepackage{amsmath,amssymb,amsthm,colortbl,xcolor,nicefrac,setspace}
\usepackage[T1]{fontenc}
\usepackage[expert,seriftt,lucidasmallscale]{lucidabr}
\usepackage{mathrsfs,hyperref,color,graphicx}
\usepackage{booktabs,array,tabularx,placeins}

% Semantic brackets, as in the local Forbes/Lucida reports.
\DeclareSymbolFont{stmry}{U}{stmry}{m}{n}
\DeclareMathDelimiter{\llbracket}{\mathopen}{stmry}{"4A}{stmry}{"71}
\DeclareMathDelimiter{\rrbracket}{\mathclose}{stmry}{"4B}{stmry}{"79}

% One-inch margins; copied from the repository's report preambles.
\setlength{\textwidth}{6.5in}
\setlength{\textheight}{9in}
\setlength{\oddsidemargin}{0mm}
\setlength{\evensidemargin}{0mm}
\setlength{\topmargin}{0mm}
\setlength{\headheight}{0mm}
\setlength{\headsep}{0mm}
\setlength{\topskip}{0mm}
\setlength{\hoffset}{0in}
\setlength{\voffset}{0in}
\newcommand{\AND}{\ensuremath{\mathbin{\&}}}
\newcommand{\IFF}{\ensuremath{\leftrightarrow}}
\newcommand{\IC}{\ensuremath{\rightarrow}}
\newcommand{\sem}[1]{\ensuremath{\llbracket #1 \rrbracket}}
\newcommand{\isa}[1]{\nolinkurl{#1}}
\onehalfspacing
\clubpenalty=10000
\widowpenalty=10000
\setlength{\emergencystretch}{3em}
\hypersetup{colorlinks=true,linkcolor=blue!50!black,
  citecolor=blue!50!black,urlcolor=blue!50!black,
  pdftitle={Why I Chose Isabelle/HOL for Autoformalizing Higher-Order Metaphysics},
  pdfauthor={},pdfsubject={Foundations, embedding architecture, and agent workflows}}

\begin{document}
\begin{center}
{\Large\bfseries Why I Chose Isabelle/HOL\par}
\vspace{4pt}
{\large Autoformalizing Higher-Order Metaphysics\par}
\vspace{8pt}
{\normalsize 7 September 2026; core status updated 9 September 2026}
\end{center}
\bigskip

\noindent I chose Isabelle/HOL because the task was to formalize the
metatheory of particular higher-order logics: their languages, proof systems,
and models, together with the relations between them. Classical higher-order
logic is a natural metalanguage for that task. Isabelle's inductive
definitions and locales support explicit representations with visible
assumptions; Sledgehammer and Nitpick support proof search and counterexample
investigation; and Isar provides a readable structure for the resulting
proofs. Those are methodological reasons for the original choice.

Lean 4 can support this work too. Its dependent types, metaprogramming, and
visible ecosystem for language-model-assisted formalization are genuine
attractions. My choice gives greater weight to Isabelle's fit for exploratory
logical metatheory and its established combination of symbolic tools.
It does not depend on the amount of code subsequently written in Isabelle,
and I have not run a controlled benchmark establishing that Isabelle agents
outperform Lean agents on this task.

The comparison requires three distinctions: the host logic versus the
represented object theory; deep versus shallow embedding; and proof
correctness versus correctness of the formalization. The main case is the
standalone Bacon--Dorr core repository. Two external comparisons,
Goodman's Purity-of-Pure question and Kirchner's implementation of Zalta's
Abstract Object Theory (AOT), illustrate the same distinctions; neither
application is included in this standalone repository. The external-tool
assessment retains its date of 7 September 2026, while the core status below
has been updated through 9 September.

\section{Foundations and the choice of representation}

Isabelle/HOL implements classical higher-order logic with simple types,
within Isabelle's generic logical framework. This makes it a congenial
metalanguage for classical reasoning about typed higher-order languages,
models, and proof systems. Its datatypes, inductive definitions, recursive
functions, and locales provide established means of expressing the relevant
structures and their assumptions~\cite{isabelle}. Lean 4 is based on dependent
type theory. Classical reasoning is available in Lean; using it does not
require recasting the source theory as constructive mathematics~\cite{lean-types}.
Both systems can support the metatheory at issue.

The closer fit of Isabelle/HOL is nevertheless a fit of \emph{metalanguage}
to task. It does not identify Bacon's logic with Isabelle's HOL. In particular,
we must not simply identify source propositions with HOL Booleans, source
identity with equivalence of truth values, or every higher-type domain with
a full function space. Those choices can settle substantive issues of
extensionality or granularity before the source axioms have been represented.
The same discipline is necessary in Lean. Neither host establishes the
adequacy of an embedding merely by checking its proofs.

In a \emph{deep embedding}, expressions are data, and typing and derivability
are independently defined relations on those data. The source-facing core uses explicit named terms, typing, declared
signatures, capture-avoiding substitution, conversion, and independently
defined deduction. Earlier de Bruijn and parametric representations remain
as supporting layers, with their source correspondences proved separately.
Semantic interpretation is a further construction.
Thus a statement of the schematic form
\[
  S\vdash_H A \quad\IFF\quad S\models_{\mathcal M} A
\]
relates two separately specified notions. Here $S$ is a theory of sentences,
$A$ is a sentence, and $\mathcal M$ is a specified model class. The typing,
language, and model-class qualifications are part of the theorem.

This architecture is useful when the research concerns substitution,
admissible rules, derivability, soundness, or completeness. It also incurs
substantial work: an agent must learn the local syntax and prove the
structural lemmas that connect its layers. Isabelle supports that work well;
it does not make it automatic. Lean could support the same extrinsic
representation. The choice is not between an Isabelle system capable of
metatheory and a Lean system incapable of it.

\section{Proof search, countermodels, and auditability}

Sledgehammer supplies an unusually mature connection between interactive
proof and external automated provers. It selects relevant facts, translates
proof obligations, and suggests proofs for replay in Isabelle. It is
particularly useful after the conceptual work has reduced a goal to a
consequence of structural lemmas. An external prover's success remains
separate from successful proof reconstruction; the accepted artifact should
be the replayed Isabelle proof~\cite{sledgehammer}.

Nitpick contributes a complementary form of investigation: searching for
counterexamples to proposed HOL statements. In a metaphysics development,
this can expose an omitted assumption in a modal inference, an overstrong
identity principle, or a proposed semantic condition that fails in a small
structure. Its finite searches and approximations must be interpreted with
their reported status. A potential counterexample requires checking, and
failure to find one establishes neither validity nor consistency. Moreover,
a deep embedding can make the relevant search difficult: a finite test of a
semantic fragment does not automatically instantiate all the closure and
schema conditions of the source theory~\cite{nitpick}.

Lean has substantial symbolic automation too. Its core \isa{grind} tactic,
for example, combines several forms of automated reasoning; Lean Copilot can
combine generated tactics with Aesop~\cite{lean-grind,copilot}. Isabelle's
advantage here is the established integration of proof search and model
finding, especially for exploratory logical metatheory. We should not
convert that practical advantage into a universal ranking of provers.

Isar also supports an auditable presentation: named assumptions, intermediate
claims, and explicit proof structure can accompany a source locator and a
notation dictionary~\cite{isar}. That is valuable when a reader must check
why a formal lemma represents a paragraph in a paper. Lean supports readable
structured proofs as well. In either system, source correspondence must be
written and maintained; it is not a consequence of choosing an attractive
surface notation.

\section{Lean's advantages for autoformalization}

The strongest current argument for Lean is the surrounding AI infrastructure.
LeanDojo established tools for extracting premises, tactics, and proof states
and for interacting with the prover. Its current repository directs new work
to LeanDojo-v2, whose proof-search interface uses Pantograph to apply tactics
and inspect resulting goals~\cite{dojo}. Lean Copilot supplies tactic
suggestion, proof search, and premise retrieval inside Lean~\cite{copilot}.
These are concrete examples of the generate--check--repair cycle that an
agent needs.

Statement autoformalization has a substantial Lean-based research literature.
ProofNet originally paired informal mathematics with \emph{Lean 3}
statements~\cite{proofnet}; it should not be described as originally a Lean 4
benchmark. Subsequent work evaluates Lean 4 translations, including corrected
ProofNet data and the research-level RLM25 benchmark~\cite{evaluation}.
Taken together, these tools and benchmarks support the judgment that Lean
currently offers a broader visible LLM/autoformalization ecosystem. They do
not tell us the composition of any proprietary model's training corpus, nor
do they establish a numerical advantage for higher-order metaphysics.

Lean's language design is a further attraction. Macros and elaborators permit
custom notation to be translated into checked terms, and tactics can be
implemented in Lean itself~\cite{lean-meta}. This makes it convenient to build
a domain language, expose proof states, and connect search to elaboration.
Isabelle has its own extensive customization facilities, including
Isabelle/ML and Eisbach~\cite{isabelle,isar}. The advantage is one of developer
workflow and available integrations, not an exclusive capability of Lean.

Dependent types also permit \emph{intrinsically typed syntax}. Schematically,
one can define a family $\mathrm{Term}(\Gamma,\tau)$, whose members are terms
of object type $\tau$ in context $\Gamma$. Application and abstraction then
have shapes such as
\[
\begin{aligned}
\mathrm{app}&:
 \mathrm{Term}(\Gamma,\sigma\IC\tau)\times
 \mathrm{Term}(\Gamma,\sigma)\longrightarrow\mathrm{Term}(\Gamma,\tau),\\
\mathrm{lam}&:
 \mathrm{Term}(\sigma::\Gamma,\tau)
 \longrightarrow\mathrm{Term}(\Gamma,\sigma\IC\tau).
\end{aligned}
\]
These are representation sketches, not executable Lean declarations.
Context-indexed variables would supply the variable case. Ill-typed
applications cannot be constructed in such a representation.

This can remove separate preservation obligations, but dependent substitution
and equality transport may become more involved. More importantly, a source
with an explicit typing relation may call for a theorem connecting that
relation to the intrinsic representation. Intrinsic typing need not hide the
source machinery: we can retain raw syntax and prove an adequacy bridge.
For this task, I preferred an explicit extrinsic representation because the
source typing judgments and their preservation properties are themselves
objects of investigation. That preference concerns what the formalization
should expose, rather than the cost of changing an existing implementation.

\section{Why project-specific embeddings change the comparison}

A model familiar with ordinary Lean mathematics can reuse conventions and
library lemmas with little adaptation. It cannot thereby know the intended
meaning of this project's \isa{Forall}, \isa{Lam}, or \isa{H_proves}, nor which
conversion or consequence relation a source theorem requires. A deep embedding in Lean
would introduce similarly project-specific definitions. General skills in
induction, tactics, and library use would still transfer, but much of the
critical knowledge would have to come from the project.

For this reason, retrieval and source fidelity are likely to dominate many
of the difficult tasks here. Useful context includes the exact statement,
its source passage, the permitted imports, typing and signature conditions,
and nearby proofs at the same logical strength. The standalone repository has a notation guide, an explanatory source
table, and an Isabelle-native dependency graph for the checked core.
These are practical foundations for an agent that retrieves the right
mathematical context. The larger parent research project keeps its
Bacon--Dorr--Goodman and AOT graphs separate.

This is a reason to expect a reduction in Lean's generic pretraining
advantage, not evidence that the advantage disappears. Research on Lean
statement autoformalization itself finds difficulty with research-level
material without appropriate context~\cite{evaluation}. A direct comparison
on previously unseen tasks from this repository would still be needed.
Existing retrieval tools would also need adaptation: Lean Copilot's documented
premise selector uses a fixed Lean/mathlib snapshot, rather than automatically
indexing any new embedded theory~\cite{copilot}.

\section{The core development and two external comparisons}

\subsection{Bacon and Bacon--Dorr: explicit logical metatheory}

Bacon's higher-order logics and the Bacon--Dorr treatment of Classicism make
the source language, inferential rules, and semantic assumptions central to
the formalization~\cite{bacon,dorr}. Isabelle is a good fit for keeping these
objects separate while proving their connections. Binding, substitution,
signatures, deduction, and model construction were therefore central
requirements when choosing a prover, rather than just infrastructure to
preserve after making the choice.

The standalone repository's September 9 checkpoint contains soundness,
sentence-set model existence, and closed strong completeness for the
Bacon--Dorr named H presentation. The book's full minimal H language also
has general-model soundness, original-signature model existence, and global
strong completeness, including open formulas and arbitrary premise sets.
The statements retain their explicit variable-stock, language, model-class,
and carrier qualifications.

For the paper's relational-type Classicism, category soundness/completeness
and single-formula action-model soundness/completeness are checked at the
stated carrier scopes. The book's full-type C development has an actual
canonical modal model and original-signature model existence for countably
declared signatures. Generic interpretation existence is now checked for
every independent full-minimal modal model: typed K/S abstraction elimination
supplies values in the chosen domains and proves the literal future-function
abstraction clause. No supplied interpreter or extra nonemptiness
assumption is added.

Generic full-type C soundness and final unrestricted modal completeness
remain open. The HOL--ZF model representation retains its explicit
foundation and documented future-restricted implication convention.
These boundaries are recorded in the standalone source guide and
verification records. This comparison note summarizes those checked
endpoints; it supplies no new proof certificate.

\subsection{Goodman: preserving the exact question}

Goodman's question asks whether Purity of Pure is consistent with the
remaining theory and a unique fundamental proposition, without assuming
Purity of Fun~\cite{goodman}. Purity of Pure says that the purity property
itself is pure; the logical stock and closure principles determine what this
requires at the relevant types. A contradiction obtained by adding Purity of
Fun would answer a stronger question. Likewise, a small structure satisfying
selected instances would not establish a model of the unrestricted theory.

The parent research project's record leaves the central question open. It separates
the canonical reconstruction of Bacon's appendix model, Goodman-specific
interpretations on those carriers, and secondary experimental models. Some
results use HOL-ZF with axiomatized ZFC and retain that relative semantic
basis. These distinctions illustrate why both proof checking and
countermodel search need source-sensitive interpretation. Keeping Goodman
in Isabelle allows it to reuse the exact background development, with
explicit bridges wherever alternative presentations are involved.

\subsection{Kirchner and Zalta: a shallow embedding with a restricted interface}

Kirchner's AOT environment uses a \emph{shallow semantic embedding}: object
expressions are represented through semantic objects in HOL, and much of the
binding infrastructure uses HOL functions. There is no general formula
datatype and independently reified derivation relation analogous to the
Bacon development. Zalta's \emph{Principia Logico-Metaphysica} supplies
the target presentation~\cite{zalta}. This substantially reduces the amount of explicit syntax
metatheory required for proving theorems within AOT.

However, the embedding is not a simple identification of AOT with HOL.
AOT includes hyperintensional relations, free logic, and the distinction
between exemplification and encoding. Kirchner supplies custom notation and
proof commands, together with an abstraction layer that restricts automation
to AOT axioms and inferential rules. This is intended to exclude theorems
that merely reflect details of the semantic implementation~\cite{kirchner}.

A Lean implementation could also exploit host binding, types, and custom
elaboration. Thus shallow embedding narrows the architectural difference
between the provers, while leaving substantial semantic and interface work.
The AOT case illustrates what a different choice of embedding can accomplish
in Isabelle; it is not the basis for the choice of a deep embedding here.
That external AOT session extends \isa{HOL-Cardinals} independently and is
not part of this standalone core repository. A shared host prover
does not supply a mathematical bridge between the theories.

\section{A concise comparison}

\noindent Table~\ref{tab:comparison} summarizes the practical tradeoffs. The
entries describe available support and this repository's needs, rather than
measured success rates.

\begin{table}[htbp]
\centering
\begin{singlespace}
\small
\renewcommand{\arraystretch}{1.18}
\begin{tabularx}{\textwidth}{@{}>{\raggedright\arraybackslash}p{1.18in}
 >{\raggedright\arraybackslash}X >{\raggedright\arraybackslash}X@{}}
\toprule
\textbf{Issue} & \textbf{Isabelle/HOL} & \textbf{Lean 4}\\
\midrule
Host foundation & Classical simple HOL; close metatheoretic fit & Dependent type theory; classical reasoning available\\
Deep embedding & Mature syntax, deduction, and model infrastructure & Equally viable; intrinsic typing is an additional option\\
Automation & Sledgehammer and integrated Nitpick exploration & Strong tactics; extensive AI integrations\\
LLM workflows & Smaller ecosystem; recent proof-state and search work & Broader visible tooling and benchmark ecosystem\\
Customization & Isabelle/ML, Eisbach, Isar, custom syntax & Macros, elaborators, and tactics written in Lean\\
Source audit & Structured Isar proofs with named assumptions & Structured proofs and source maps can be built\\
Shallow embedding & External AOT implementation illustrates the approach & Host binding and elaboration also support the approach\\
Fit for this task & Classical metatheory with integrated proof and counterexample tools & Viable metatheory with dependent types and extensive AI tooling\\
\bottomrule
\end{tabularx}
\end{singlespace}
\caption{Tradeoffs for autoformalizing higher-order metaphysics.}
\label{tab:comparison}
\end{table}
\FloatBarrier

\section{What the choice of tool does not settle}

The special difficulty in autoformalizing philosophy is that the intended
formal statement may itself be at issue. A proof assistant can accept a
proof of a perfectly well-formed statement that misrepresents the source.
This occurs in mathematics too, as recent work on evaluating statement
autoformalization emphasizes~\cite{evaluation}; philosophy adds sustained
questions about the interpretation of the primitives, axioms, and semantics.
In this setting, a higher proof-success rate is useful only when it measures
success on the right formalization.

For example, identifying necessarily equivalent propositions when the target
theory distinguishes their identity can erase the intended hyperintensional
structure. Restricting a quantifier to an unintended stock can make a schema
easier to prove. Importing a stronger calculus can turn an open derivability
question into an immediate consequence. A kernel checks the resulting
formal argument; it does not compare these choices with the philosophical
text. An effective workflow must therefore audit the translation as well as
the proof.

The resulting workflow should keep four obligations explicit:
\begin{enumerate}
\item \textbf{Fix the target.} Record the source passage, formal statement,
logical strength, signature, and semantic scope before generating a proof.
Have the source correspondence reviewed separately from tactic success.
\item \textbf{Retrieve within that scope.} Use the notation dictionary,
source matrix, theorem statements, and the appropriate Isabelle-native graph.
Keep the independent theory families separate and expose the assumptions of
candidate premises.
\item \textbf{Check and repair incrementally.} Give the agent structured proof
states and errors; combine short proof attempts with Sledgehammer and, where
suitable, Nitpick. Retain replayable proofs and validate candidate models.
Serialize Isabelle builds in accordance with the repository's existing rules.
\item \textbf{Evaluate on faithful results.} Use held-out source passages and
lemmas from the repository. Record translation accuracy, completed checked
proofs, premise discipline, human corrections, and time or cost. Do not count
an easier substituted statement as a successful formalization.
\end{enumerate}

There is evidence that such a pilot is feasible. The IsaMini line of work
studies a simplified proof language combined with Isabelle automation;
Stepwise exposes fine-grained Isabelle proof states through a REPL and
combines learned search with symbolic repair~\cite{isamini,stepwise}. These
are research demonstrations on their own tasks. They do not establish
compatibility with the local toolchain or success on metaphysical source
translation, which the pilot would need to test.

My reason for choosing Isabelle was the combination of a suitable
metalanguage, mature proof and counterexample tools, and a proof format
that makes assumptions and intermediate reasoning inspectable. These are
useful features before any project-specific development exists. Lean's
advantages remain relevant, especially when a project's main objective is
to use the broadest available AI formalization infrastructure. Neither
choice replaces the need to check that the formalization expresses the
intended theory. That requirement is central to the project, whichever
prover one uses.

\medskip
\begin{singlespace}
\small
\noindent\textbf{Repository basis.} The core status was refreshed from
\isa{STATUS.md}, \isa{docs/SOURCE_CORRESPONDENCE.md}, and
\isa{docs/VERIFICATION.md} at the standalone repository's checked
9 September 2026 checkpoint \isa{41bb91f}. The main H and C scopes and
generic interpretation-existence audit are recorded there. The original
comparison's external-tool assessment and bibliography date from
7 September; no new comparative benchmark or external-agent compatibility
test was run for this revision. The parent report is retained unchanged.
\par\smallskip
\url{https://github.com/fitelson/bacon-dorr-isabelle}
\end{singlespace}

\bigskip
\begin{singlespace}
\footnotesize
\raggedright
\begin{thebibliography}{99}
\bibitem{isabelle} Isabelle. \emph{Documentation for Isabelle2025-2}.
\url{https://isabelle.in.tum.de/documentation.html}.
\bibitem{lean-types} Lean. \emph{The Lean Language Reference}, ``The Type
System'' and ``Axioms.''
\url{https://lean-lang.org/doc/reference/latest/}.
\bibitem{sledgehammer} Jasmin Blanchette, with contributions from Martin
Desharnais, Lawrence C.\ Paulson, and Lukas Bartl (2026).
\emph{Hammering Away: A User's Guide to Sledgehammer for Isabelle/HOL}.
\url{https://isabelle.in.tum.de/doc/sledgehammer.pdf}.
\bibitem{nitpick} Jasmin Blanchette (2026).
\emph{Picking Nits: A User's Guide to Nitpick for Isabelle/HOL}.
\url{https://isabelle.in.tum.de/doc/nitpick.pdf}.
\bibitem{lean-grind} Lean. \emph{The Lean Language Reference}, ``The
\isa{grind} Tactic.''
\url{https://lean-lang.org/doc/reference/latest/The--grind--tactic/}.
\bibitem{copilot} LeanDojo contributors. \emph{Lean Copilot}: project and
usage documentation. \url{https://github.com/lean-dojo/LeanCopilot}.
\bibitem{isar} Makarius Wenzel. \emph{The Isabelle/Isar Reference Manual}.
\url{https://isabelle.in.tum.de/doc/isar-ref.pdf}.
\bibitem{dojo} LeanDojo contributors. \emph{LeanDojo} and
\emph{LeanDojo-v2}: project documentation.
\url{https://github.com/lean-dojo/LeanDojo};
\url{https://github.com/lean-dojo/LeanDojo-v2}.
\bibitem{proofnet} Zhangir Azerbayev et al. (2023).
``ProofNet: Autoformalizing and Formally Proving Undergraduate-Level
Mathematics.'' \url{https://arxiv.org/abs/2302.12433}.
\bibitem{evaluation} Auguste Poiroux, Gail Weiss, Viktor Kun\v{c}ak, and
Antoine Bosselut (2025). ``Reliable Evaluation and Benchmarks for Statement
Autoformalization.'' \emph{EMNLP 2025}, pp.~17947--17969.
\url{https://aclanthology.org/2025.emnlp-main.907/}.
\bibitem{lean-meta} Lean. \emph{The Lean Language Reference}, ``Elaboration
and Compilation,'' ``Tactic Proofs,'' and ``Notations and Macros.''
\url{https://lean-lang.org/doc/reference/latest/}.
\bibitem{bacon} Andrew Bacon (2024).
\emph{A Philosophical Introduction to Higher-Order Logics}. Routledge.
\url{https://doi.org/10.4324/9781003039181}.
\bibitem{dorr} Andrew Bacon and Cian Dorr (2022). ``Classicism.'' Draft,
1 July 2022. \url{https://andrew-bacon.github.io/papers/Classicism.pdf}.
Numbering follows the draft specified in the repository's source guide.
\bibitem{goodman} Jeremy Goodman (2026). ``Is Purity of Pure consistent with
Bacon's Logical Combinatorialism?'' Project notes, July 2026.
Unpublished source retained in the parent research project, not redistributed here.
\bibitem{zalta} Edward N.\ Zalta (2026). \emph{Principia Logico-Metaphysica}.
Draft/excerpt, 15 May 2026, with critical theoretical contributions by
Daniel Kirchner and Uri Nodelman. Source copy retained in the parent
research project, not redistributed here.
\bibitem{kirchner} Daniel Kirchner (2024). ``Computer-Verified Foundations
of Metaphysics.'' \emph{KI -- K\"unstliche Intelligenz} 38:95--98.
\url{https://doi.org/10.1007/s13218-024-00834-z}.
The maintained implementation is at \url{https://github.com/ekpyron/AOT};
it is independent of the standalone Bacon--Dorr core.
\bibitem{isamini} Qiyuan Xu, Renxi Wang, Peixin Wang, Haonan Li, and Conrad
Watt. ``A Minimalist Proof Language for Neural Theorem Proving over
Isabelle/HOL.'' arXiv:2507.18885, initially submitted 2025; current revision.
\url{https://arxiv.org/abs/2507.18885}.
\bibitem{stepwise} Baoding He et al. (2026). ``Stepwise: Neuro-Symbolic
Proof Search for Automated Systems Verification.'' arXiv:2603.19715.
\url{https://arxiv.org/abs/2603.19715}.
\end{thebibliography}
\end{singlespace}
\end{document}
